\begin{CNbox}[unbreakable]{obj}{Learning Objectives}

By the end of this module, you will be able to:

\begin{itemize}
\tightlist
\item
  Explain WHY mobility management is essential in cellular networks
\item
  Distinguish between Location Areas, Routing Areas, and Tracking Areas
\item
  Describe cell reselection mechanisms in idle mode
\item
  Classify and compare handover types across generations
\item
  Analyze the handover decision process using measurement events
\item
  Trace the LTE X2 handover procedure step by step
\item
  Identify handover failures and their mitigation strategies
\item
  Compare mobility evolution from 2G through 5G
\end{itemize}

\end{CNbox}

\Needspace{14\baselineskip}

\CNsection{1}{Why Mobility Management Is Needed}

\CNchained\CNsubsection{}{WHY}

Users are mobile — they walk, drive, and travel while expecting uninterrupted voice calls, video streams, and data sessions. Without mobility management, a user crossing a cell boundary would experience immediate service disconnection.

\CNsubsection{}{WHAT}

Mobility management is the set of network procedures that:

\begin{itemize}
\tightlist
\item
  \textbf{Track} the user’s location (so the network can reach them for incoming calls/\allowbreak{}data)
\item
  \textbf{Maintain} service continuity when the user moves between cells (handover)
\item
  \textbf{Optimize} radio resource usage by connecting users to the best-serving cell
\end{itemize}

\Needspace{14\baselineskip}

\CNsubsection{}{HOW}

The network achieves this through two complementary mechanisms:

{\def\LTcaptype{none} % do not increment counter
\Needspace{8\baselineskip}
\begin{longtable}[c]{@{}
  >{\raggedright\arraybackslash\hspace{0pt}}m{(\linewidth - 4\tabcolsep) * \real{0.2593}}
  >{\raggedright\arraybackslash\hspace{0pt}}m{(\linewidth - 4\tabcolsep) * \real{0.4074}}
  >{\raggedright\arraybackslash\hspace{0pt}}m{(\linewidth - 4\tabcolsep) * \real{0.3333}}@{}}
\toprule\noalign{}
\rowcolor{CNink}\CNth{State} & \CNth{Mechanism} & \CNth{Purpose} \\
\CNheadrulesep
\endhead
\bottomrule\noalign{}
\endlastfoot
\textbf{Idle mode} & Cell reselection + location/\allowbreak{}tracking updates & Network knows user’s approximate location \\
\textbf{Connected mode} & Handover & Seamless transfer of active connection \\
\end{longtable}
}

\CNkeepfig{m04-lu-tradeoff}{3}

\CNsubsection{}{The Fundamental Trade-off}

\CNfigure{m04-lu-tradeoff}{The location-update trade-off: paging accuracy versus signalling load}

The network must balance:

\begin{itemize}
\tightlist
\item
  \textbf{Signaling load} (updates consume radio resources)
\item
  \textbf{Paging load} (larger areas = more cells to page)
\item
  \textbf{Service latency} (time to reach the user for incoming services)
\end{itemize}

\Needspace{14\baselineskip}

\CNsection{2}{Location Management: LA vs RA vs TA}

\CNchained\CNsubsection{}{WHY}

The network needs to know WHERE a user is (approximately) to deliver incoming calls or data. Different generations use different grouping strategies to balance signaling vs. paging efficiency.

\Needspace{15\baselineskip}

\CNsubsection{}{WHAT}

{\def\LTcaptype{none} % do not increment counter
\Needspace{11\baselineskip}
\begin{longtable}[c]{@{}
  >{\raggedright\arraybackslash\hspace{0pt}}m{(\linewidth - 6\tabcolsep) * \real{0.2292}}
  >{\raggedright\arraybackslash\hspace{0pt}}m{(\linewidth - 6\tabcolsep) * \real{0.2708}}
  >{\raggedright\arraybackslash\hspace{0pt}}m{(\linewidth - 6\tabcolsep) * \real{0.2083}}
  >{\raggedright\arraybackslash\hspace{0pt}}m{(\linewidth - 6\tabcolsep) * \real{0.2917}}@{}}
\toprule\noalign{}
\rowcolor{CNink}\CNth{Generation} & \CNth{Area Concept} & \CNth{Used For} & \CNth{Typical Size} \\
\CNheadrulesep
\endhead
\bottomrule\noalign{}
\endlastfoot
2G (GSM) & \textbf{Location Area (LA)} & Circuit-switched services & 20–100 cells \\
2.5G (GPRS) & \textbf{Routing Area (RA)} & Packet-switched services & Subset of LA, 10–50 cells \\
4G (LTE) & \textbf{Tracking Area (TA)} & All services (all-IP) & 30–80 cells \\
5G (NR) & \textbf{Tracking Area} + RNA & All services + RAN-level & Configurable per UE \\
\end{longtable}
}

\CNkeepfig{m04-la-cells}{3}

\CNsubsection{}{HOW — Location Area (2G)}

\CNfigure{m04-la-cells}{A Location Area groups many cells}

\begin{itemize}
\tightlist
\item
  UE performs \textbf{Location Area Update (LAU)} when crossing LA boundary
\item
  Network pages ALL cells in the LA for incoming call
\item
  Managed by MSC/\allowbreak{}VLR
\end{itemize}

\CNsubsection{}{HOW — Routing Area (GPRS)}

\begin{itemize}
\tightlist
\item
  RA is a \textbf{subset} of an LA (one LA contains multiple RAs)
\item
  UE performs \textbf{Routing Area Update (RAU)} for packet services
\item
  Managed by SGSN
\item
  Allows finer-grained location for data users
\end{itemize}

\CNsubsection{}{HOW — Tracking Area (LTE/\allowbreak{}5G)}

\begin{itemize}
\tightlist
\item
  \textbf{Tracking Area List (TAL)}: UE is assigned a LIST of TAs — no update needed when moving within the list
\item
  Reduces signaling for high-mobility users (larger TAL)
\item
  Low-mobility users get smaller TAL to reduce paging area
\end{itemize}

\tcbset{CNcodemode/.style={unbreakable}}\def\CNcodelang{}

\begin{CNplaincode}
\begin{Verbatim}
UE TAL = {TA1, TA2, TA3, TA4}

Moving TA1→TA2→TA3→TA4: NO update needed
Moving TA4→TA5 (outside TAL): Tracking Area Update triggered
\end{Verbatim}
\end{CNplaincode}

\CNsubsection{}{5G Enhancement: RAN-based Notification Area (RNA)}

\begin{itemize}
\tightlist
\item
  In \textbf{RRC\_\allowbreak{}INACTIVE} state (new in 5G), the gNB manages a smaller RNA
\item
  Reduces core network signaling
\item
  Two-level mobility: RNA (RAN-managed) + TA (core-managed)
\end{itemize}

\Needspace{14\baselineskip}

\CNsection{3}{Cell Reselection (Idle Mode Mobility)}

\CNchained\CNsubsection{}{WHY}

Even when not in an active call/\allowbreak{}session, the UE must camp on the best cell to:

\begin{itemize}
\tightlist
\item
  Receive paging messages efficiently
\item
  Ensure good signal quality when a connection is needed
\item
  Avoid unnecessary location updates
\end{itemize}

\CNsubsection{}{WHAT}

Cell reselection is the \textbf{UE-autonomous} process of selecting a new cell to camp on while in idle mode. The network provides parameters via System Information Broadcasts (SIBs), but the UE makes the decision independently.

\Needspace{11\baselineskip}

\CNsubsection{}{HOW — Cell Reselection Criteria}

\CNchained\CNsubsubsection{Ranking Criterion (S-criterion)}

A cell is suitable for camping if it meets the S-criterion:

\CNdisplay{
\begin{aligned}
S_{\mathrm{rxlev}} &= Q_{\mathrm{rxlevmeas}} - \left(Q_{\mathrm{rxlevmin}} + Q_{\mathrm{rxlevminoffset}}\right) - P_{\mathrm{compensation}} > 0\\
S_{\mathrm{qual}} &= Q_{\mathrm{qualmeas}} - \left(Q_{\mathrm{qualmin}} + Q_{\mathrm{qualminoffset}}\right) > 0
\end{aligned}
}

Where:

\begin{itemize}
\tightlist
\item
  \CNcode{Qrxlevmeas}: Measured RSRP (dBm)
\item
  \CNcode{Qrxlevmin}: Minimum required RSRP (-124 to -44~dBm, typical: -124~dBm)
\item
  \CNcode{Qqualmeas}: Measured RSRQ (dB)
\item
  \CNcode{Qqualmin}: Minimum required RSRQ (-20 to -3~dB, typical: -20~dB)
\end{itemize}

\CNsubsubsection{Cell Ranking (R-criterion)}

For intra-frequency reselection, cells are ranked:

\CNdisplay{
\begin{aligned}
R_s &= Q_{\mathrm{meas},s} + Q_{\mathrm{Hyst}} &&\text{(serving cell)}\\
R_n &= Q_{\mathrm{meas},n} - Q_{\mathrm{offset}} &&\text{(neighbour cell)}
\end{aligned}
}

Reselect to the neighbour if $R_n > R_s$ holds for the duration $T_{\mathrm{reselection}}$.

\Needspace{12\baselineskip}

\CNsubsubsection{Priority-Based Reselection (Inter-frequency/\allowbreak{}Inter-RAT)}

{\def\LTcaptype{none} % do not increment counter
\Needspace{9\baselineskip}
\begin{longtable}[c]{@{}
  >{\raggedright\arraybackslash\hspace{0pt}}m{(\linewidth - 2\tabcolsep) * \real{0.5556}}
  >{\raggedright\arraybackslash\hspace{0pt}}m{(\linewidth - 2\tabcolsep) * \real{0.4444}}@{}}
\toprule\noalign{}
\rowcolor{CNink}\CNth{Priority} & \CNth{Action} \\
\CNheadrulesep
\endhead
\bottomrule\noalign{}
\endlastfoot
Higher priority layer & Reselect if neighbor meets ThreshX,High \\
Equal priority layer & Use ranking (R-criterion) \\
Lower priority layer & Reselect only if serving \textless{} ThreshServing,Low AND neighbor \textgreater{} ThreshX,Low \\
\end{longtable}
}

\CNsubsection{}{HOW — Hysteresis and Timers}

\textbf{Hysteresis (QHyst)} prevents ping-pong reselection:

\begin{itemize}
\tightlist
\item
  Typical values: 2–6~dB
\item
  Adds ``stickiness'' to the serving cell
\end{itemize}

\textbf{Treselection timer} (typically 1–4 seconds):

\begin{itemize}
\tightlist
\item
  Neighbor must remain better for this duration before reselection occurs
\item
  Prevents reselection due to momentary fading
\end{itemize}

\CNfigure{m04-reselection}{Cell reselection: the neighbour must stay above serving + QHyst for Treselection}

\Needspace{18\baselineskip}

\CNsubsection{}{Typical Cell Reselection Parameters}

{\def\LTcaptype{none} % do not increment counter
\Needspace{14\baselineskip}
\begin{longtable}[c]{@{}
  >{\raggedright\arraybackslash\hspace{0pt}}m{(\linewidth - 4\tabcolsep) * \real{0.3235}}
  >{\raggedright\arraybackslash\hspace{0pt}}m{(\linewidth - 4\tabcolsep) * \real{0.4118}}
  >{\raggedright\arraybackslash\hspace{0pt}}m{(\linewidth - 4\tabcolsep) * \real{0.2647}}@{}}
\toprule\noalign{}
\rowcolor{CNink}\CNth{Parameter} & \CNth{Typical Value} & \CNth{Purpose} \\
\CNheadrulesep
\endhead
\bottomrule\noalign{}
\endlastfoot
QHyst & 4~dB & Serving cell hysteresis \\
Treselection & 2~s & Stability timer \\
Qrxlevmin & -124~dBm & Minimum RSRP for cell selection \\
Qqualmin & -20~dB & Minimum RSRQ for cell selection \\
ThreshServing,Low & 4 (= 4~dB) & Threshold to leave serving (to lower priority) \\
ThreshX,High & 8 (= 8~dB) & Threshold to enter higher priority layer \\
\end{longtable}
}

\Needspace{14\baselineskip}

\CNsection{4}{Handover Types}

\CNchained\CNsubsection{}{WHY}

Different mobility scenarios require different handover strategies. A user walking between co-channel cells needs a different mechanism than a user on a high-speed train crossing between LTE and 3G coverage areas.

\CNkeepfig{m04-ho-types}{3}

\CNsubsection{}{WHAT — Classification Overview}

\CNfigure{m04-ho-types}{Classification of handovers}

\Needspace{11\baselineskip}

\CNsubsection{4.1}{Intra-Frequency Handover}

\CNchained\CNsubsubsection{WHY}

Most common handover — occurs when user moves between cells operating on the \textbf{same carrier frequency}. This is the typical scenario in dense urban deployments.

\CNsubsubsection{WHAT}

\begin{itemize}
\tightlist
\item
  Source and target cells use the same frequency (e.g., both on 1800~MHz Band 3)
\item
  UE can measure neighbors WITHOUT measurement gaps
\item
  Fastest and simplest handover type
\end{itemize}

\CNsubsubsection{HOW}

\begin{itemize}
\tightlist
\item
  UE continuously measures RSRP/\allowbreak{}RSRQ of serving and intra-frequency neighbors
\item
  No interruption to measurements needed (same frequency = can measure while receiving)
\item
  Triggered by Event A3: Neighbor becomes offset better than serving
\end{itemize}

\textbf{Typical scenario}: Walking between adjacent cells in a shopping mall (all cells on same frequency)

\Needspace{11\baselineskip}

\CNsubsection{4.2}{Inter-Frequency Handover}

\CNchained\CNsubsubsection{WHY}

Networks deploy multiple frequency layers for capacity (e.g., macro on 800~MHz for coverage, small cells on 2600~MHz for capacity). Users must move between these layers.

\CNsubsubsection{WHAT}

\begin{itemize}
\tightlist
\item
  Source and target cells operate on \textbf{different carrier frequencies}
\item
  UE needs \textbf{measurement gaps} to tune its receiver to other frequencies
\item
  Slightly longer interruption than intra-frequency
\end{itemize}

\CNsubsubsection{HOW}

\begin{itemize}
\tightlist
\item
  Network configures measurement gaps (6~ms every 40 or 80~ms)
\item
  UE measures inter-frequency neighbors during gaps
\item
  Triggered by Event A4 (neighbor on other frequency becomes better than threshold) or A5 (serving degrades AND neighbor is good)
\end{itemize}

\textbf{Typical scenario}: User moves from indoor small cell (2600~MHz) to outdoor macro (800~MHz)

\Needspace{11\baselineskip}

\CNsubsection{4.3}{Inter-RAT Handover}

\CNchained\CNsubsubsection{WHY}

Not all areas have LTE/\allowbreak{}5G coverage. When LTE signal degrades, the UE may need to fall back to 3G (UMTS) or even 2G (GSM) to maintain service. Also used for CS Fallback (voice calls).

\CNsubsubsection{WHAT}

\begin{itemize}
\tightlist
\item
  Handover between different radio access technologies
\item
  Examples: LTE \ensuremath{\to} UMTS, LTE \ensuremath{\to} GSM, NR \ensuremath{\to} LTE
\item
  Most complex type — different protocols, different core networks
\end{itemize}

\CNsubsubsection{HOW}

\begin{itemize}
\tightlist
\item
  UE measures inter-RAT neighbors during measurement gaps
\item
  Triggered by Event B1 (inter-RAT neighbor above threshold) or B2 (serving below threshold AND inter-RAT above threshold)
\item
  Requires inter-RAT coordination between eNB and RNC/\allowbreak{}BSC
\end{itemize}

\textbf{Typical scenario}:

\begin{itemize}
\tightlist
\item
  VoLTE not available \ensuremath{\to} CS Fallback to 3G for voice call
\item
  User drives into rural area with only 3G coverage
\end{itemize}

\Needspace{11\baselineskip}

\CNsubsection{4.4}{Hard Handover (Break-Before-Make)}

\CNchained\CNsubsubsection{WHY}

When simultaneous connection to two cells is not possible (different frequencies, different technologies, or system design choice), the old connection must be released before the new one is established.

\CNsubsubsection{WHAT}

\begin{itemize}
\tightlist
\item
  Connection to source cell is \textbf{released} before connecting to target
\item
  Brief interruption in service (typically 20–50~ms in LTE)
\item
  Used in: GSM, LTE, 5G NR
\end{itemize}

\CNkeepfig{m04-hard-ho-gap}{2}

\CNsubsubsection{HOW}

\CNfigure{m04-hard-ho-gap}{Hard handover: break before make}

\begin{itemize}
\tightlist
\item
  \textbf{Advantage}: Simpler, no need for combining/\allowbreak{}splitting
\item
  \textbf{Disadvantage}: Brief interruption (acceptable for data, noticeable for voice in 2G)
\end{itemize}

\Needspace{11\baselineskip}

\CNsubsection{4.5}{Soft Handover (Make-Before-Break) — 3G UMTS Specific}

\CNchained\CNsubsubsection{WHY}

CDMA-based 3G systems can receive the same signal from multiple cells simultaneously (same frequency, code division). This enables seamless handover with NO interruption.

\CNsubsubsection{WHAT}

\begin{itemize}
\tightlist
\item
  UE is connected to \textbf{multiple cells simultaneously} (Active Set)
\item
  Signals from multiple cells are \textbf{combined} for diversity gain
\item
  Unique to CDMA/\allowbreak{}WCDMA (3G) — NOT used in LTE or 5G
\end{itemize}

\CNsubsubsection{HOW — Active Set Management}

\tcbset{CNcodemode/.style={unbreakable}}\def\CNcodelang{}

\begin{CNplaincode}
\begin{Verbatim}
Active Set Events (3GPP UMTS):
- Event 1A: Add cell to Active Set (neighbor strong enough)
- Event 1B: Remove cell from Active Set (cell too weak)
- Event 1C: Replace cell in Active Set (better candidate found)
\end{Verbatim}
\end{CNplaincode}

\CNfigure{m04-soft-ho-region}{Soft handover region: the UE is connected to both cells}

\begin{itemize}
\tightlist
\item
  \textbf{Active Set}: Set of cells currently serving the UE (max 3–6 cells)
\item
  \textbf{Soft handover}: Between cells of same Node B or different Node Bs, same frequency
\item
  \textbf{Softer handover}: Between sectors of the SAME Node B (MRC combining at Node B)
\end{itemize}

\textbf{Advantages}:

\begin{itemize}
\tightlist
\item
  No interruption
\item
  Macro-diversity gain (2–3~dB)
\item
  Reduced ping-pong
\end{itemize}

\textbf{Disadvantages}:

\begin{itemize}
\tightlist
\item
  Uses resources in multiple cells simultaneously
\item
  Increased backhaul and processing requirements
\item
  Adds complexity to power control
\end{itemize}

\Needspace{14\baselineskip}

\CNsection{5}{Handover Decision Process}

\CNchained\CNsubsection{}{WHY}

The network must decide WHEN to handover, to WHICH cell, and ensure the decision is robust against fading and noise. A poorly timed handover causes drops; a well-tuned one is imperceptible to the user.

\CNsubsection{}{WHAT}

The handover decision in LTE is based on:

\begin{enumerate}
\def\labelenumi{\arabic{enumi}.}
\tightlist
\item
  \textbf{UE Measurement Reports} — periodic or event-triggered
\item
  \textbf{Measurement Events} — conditions that trigger the report
\item
  \textbf{Hysteresis and Time-to-Trigger} — filtering to avoid false triggers
\item
  \textbf{Network decision} — eNB evaluates report and decides
\end{enumerate}

\Needspace{11\baselineskip}

\Needspace{15\baselineskip}

\CNsubsection{5.1}{Measurement Reports (RSRP and RSRQ)}

\CNchained\CNsubsubsection{WHAT}

{\def\LTcaptype{none} % do not increment counter
\Needspace{9\baselineskip}
\begin{longtable}[c]{@{}
  >{\raggedright\arraybackslash\hspace{0pt}}m{(\linewidth - 6\tabcolsep) * \real{0.2286}}
  >{\raggedright\arraybackslash\hspace{0pt}}m{(\linewidth - 6\tabcolsep) * \real{0.3143}}
  >{\raggedright\arraybackslash\hspace{0pt}}m{(\linewidth - 6\tabcolsep) * \real{0.2000}}
  >{\raggedright\arraybackslash\hspace{0pt}}m{(\linewidth - 6\tabcolsep) * \real{0.2571}}@{}}
\toprule\noalign{}
\rowcolor{CNink}\CNth{Metric} & \CNth{Full Name} & \CNth{Range} & \CNth{Purpose} \\
\CNheadrulesep
\endhead
\bottomrule\noalign{}
\endlastfoot
\textbf{RSRP} & Reference Signal Received Power & -140 to -44~dBm & Signal strength (coverage) \\
\textbf{RSRQ} & Reference Signal Received Quality & -20 to -3~dB & Signal quality (accounts for interference + noise) \\
\textbf{SINR} & Signal-to-Interference-plus-Noise Ratio & -23 to 40~dB & Overall link quality \\
\end{longtable}
}

\CNsubsubsection{HOW}

\CNdisplay{\mathrm{RSRP} = \text{power of one resource element carrying the reference signal}}

\CNdisplay{\mathrm{RSRQ} = \frac{N \times \mathrm{RSRP}}{\mathrm{RSSI}}}

where $N$ is the number of RBs and RSSI is the total received wideband power.

\begin{itemize}
\tightlist
\item
  RSRP answers “how strong is \textbf{my} cell’s signal?”
\item
  RSRQ answers “how good is \textbf{my} cell’s signal relative to everything else I hear?”
\end{itemize}

\textbf{When to use which}:

\begin{itemize}
\tightlist
\item
  RSRP-based handover: Good for coverage-limited scenarios (rural)
\item
  RSRQ-based handover: Good for interference-limited scenarios (dense urban)
\end{itemize}

\Needspace{11\baselineskip}

\Needspace{20\baselineskip}

\CNsubsection{5.2}{LTE Measurement Events (A1–A6)}

\CNchained\CNsubsubsection{WHAT}

{\def\LTcaptype{none} % do not increment counter
\Needspace{14\baselineskip}
\begin{longtable}[c]{@{}
  >{\raggedright\arraybackslash\hspace{0pt}}m{(\linewidth - 4\tabcolsep) * \real{0.2500}}
  >{\raggedright\arraybackslash\hspace{0pt}}m{(\linewidth - 4\tabcolsep) * \real{0.3929}}
  >{\raggedright\arraybackslash\hspace{0pt}}m{(\linewidth - 4\tabcolsep) * \real{0.3571}}@{}}
\toprule\noalign{}
\rowcolor{CNink}\CNth{Event} & \CNth{Condition} & \CNth{Use Case} \\
\CNheadrulesep
\endhead
\bottomrule\noalign{}
\endlastfoot
\textbf{A1} & Serving \textgreater{} Threshold & Stop inter-freq measurements (serving is good enough) \\
\textbf{A2} & Serving \textless{} Threshold & Start inter-freq/\allowbreak{}inter-RAT measurements \\
\textbf{A3} & Neighbor \textgreater{} Serving + Offset & \textbf{Primary intra-freq handover trigger} \\
\textbf{A4} & Neighbor \textgreater{} Threshold & Inter-frequency handover \\
\textbf{A5} & Serving \textless{} Thresh1 AND Neighbor \textgreater{} Thresh2 & Inter-freq/\allowbreak{}inter-RAT handover \\
\textbf{A6} & Neighbor \textgreater{} SCell + Offset & Secondary cell (CA) change \\
\end{longtable}
}

\Needspace{11\baselineskip}

\CNsubsubsection{Event B (Inter-RAT):}

{\def\LTcaptype{none} % do not increment counter
\Needspace{8\baselineskip}
\begin{longtable}[c]{@{}
  >{\raggedright\arraybackslash\hspace{0pt}}m{(\linewidth - 4\tabcolsep) * \real{0.2500}}
  >{\raggedright\arraybackslash\hspace{0pt}}m{(\linewidth - 4\tabcolsep) * \real{0.3929}}
  >{\raggedright\arraybackslash\hspace{0pt}}m{(\linewidth - 4\tabcolsep) * \real{0.3571}}@{}}
\toprule\noalign{}
\rowcolor{CNink}\CNth{Event} & \CNth{Condition} & \CNth{Use Case} \\
\CNheadrulesep
\endhead
\bottomrule\noalign{}
\endlastfoot
\textbf{B1} & Inter-RAT neighbor \textgreater{} Threshold & Redirect to other RAT \\
\textbf{B2} & Serving \textless{} Thresh1 AND Inter-RAT \textgreater{} Thresh2 & Fallback (e.g., LTE\ensuremath{\to}3G) \\
\end{longtable}
}

\Needspace{23\baselineskip}

\CNsubsubsection{HOW — Event A3 in Detail (Most Important)}

Event A3 entering condition:\CNdisplay{\mathrm{Mn} + \mathrm{Ofn} + \mathrm{Ocn} - \mathrm{Hys} > \mathrm{Ms} + \mathrm{Ofs} + \mathrm{Ocs} + \mathrm{Off}}

{\def\LTcaptype{none} % do not increment counter
\Needspace{18\baselineskip}
\begin{longtable}[c]{@{}cc@{}}
\toprule\noalign{}
\rowcolor{CNink}\CNth{Symbol} & \CNth{Meaning} \\
\CNheadrulesep
\endhead
\bottomrule\noalign{}
\endlastfoot
Mn & Measurement of the neighbour cell (RSRP or RSRQ) \\
Ms & Measurement of the serving cell \\
Ofn & Frequency-specific offset for the neighbour \\
Ofs & Frequency-specific offset for the serving cell \\
Ocn & Cell-specific offset for the neighbour \\
Ocs & Cell-specific offset for the serving cell \\
Hys & Hysteresis (positive value) \\
Off & A3 offset (the main tunable parameter) \\
\end{longtable}
}

Simplified (same frequency, no cell-specific offsets):\CNdisplay{\mathrm{RSRP}_{\mathrm{neighbour}} - \mathrm{Hysteresis} > \mathrm{RSRP}_{\mathrm{serving}} + \mathrm{A3\_Offset}}

\Needspace{11\baselineskip}

\CNsubsection{5.3}{Hysteresis and Time-to-Trigger}

\CNchained\CNsubsubsection{WHY}

Radio signals fluctuate due to fading. Without filtering, the network would trigger handovers on every temporary fade, causing excessive signaling and ping-pong.

\Needspace{13\baselineskip}

\CNsubsubsection{WHAT}

Two mechanisms prevent false triggering:

{\def\LTcaptype{none} % do not increment counter
\Needspace{8\baselineskip}
\begin{longtable}[c]{@{}
  >{\raggedright\arraybackslash\hspace{0pt}}m{(\linewidth - 4\tabcolsep) * \real{0.3056}}
  >{\raggedright\arraybackslash\hspace{0pt}}m{(\linewidth - 4\tabcolsep) * \real{0.2778}}
  >{\raggedright\arraybackslash\hspace{0pt}}m{(\linewidth - 4\tabcolsep) * \real{0.4167}}@{}}
\toprule\noalign{}
\rowcolor{CNink}\CNth{Parameter} & \CNth{Function} & \CNth{Typical Value} \\
\CNheadrulesep
\endhead
\bottomrule\noalign{}
\endlastfoot
\textbf{Hysteresis (Hys)} & Signal must exceed threshold by this margin & 1–3~dB \\
\textbf{Time-to-Trigger (TTT)} & Condition must be sustained for this duration & 40–640~ms \\
\end{longtable}
}

\CNkeepfig{m04-a3-visual}{2}

\CNsubsubsection{HOW — Visualization}

\CNfigure{m04-a3-visual}{Event A3: entry threshold, hysteresis and time-to-trigger}

\textbf{Combined effect}:

\begin{itemize}
\tightlist
\item
  Hysteresis provides amplitude filtering (dB domain)
\item
  TTT provides time filtering (time domain)
\item
  Together they provide robust handover triggering
\end{itemize}

\Needspace{16\baselineskip}

\CNsubsubsection{Engineering Trade-off}

{\def\LTcaptype{none} % do not increment counter
\Needspace{13\baselineskip}
\begin{longtable}[c]{@{}
  >{\raggedright\arraybackslash\hspace{0pt}}m{(\linewidth - 4\tabcolsep) * \real{0.3333}}
  >{\raggedright\arraybackslash\hspace{0pt}}m{(\linewidth - 4\tabcolsep) * \real{0.3333}}
  >{\raggedright\arraybackslash\hspace{0pt}}m{(\linewidth - 4\tabcolsep) * \real{0.3333}}@{}}
\toprule\noalign{}
\rowcolor{CNink} & \CNth{Low Hys + Short TTT} & \CNth{High Hys + Long TTT} \\
\CNheadrulesep
\endhead
\bottomrule\noalign{}
\endlastfoot
Handover speed & Fast handover & Slow handover \\
Main risk & Ping-pong & Too-late HO \\
Good for & Slow users & Fast users \\
Signalling & More signalling & Less signalling \\
Typical setting & Urban: Hys = 2~dB, TTT = 160~ms & Highway: Hys = 3~dB, TTT = 80~ms \\
\end{longtable}
}

\Needspace{11\baselineskip}

\CNsubsection{5.4}{Handover Margin}

\CNchained\CNsubsubsection{WHAT}

The \textbf{total handover margin} is the combined effect of A3 offset, hysteresis, and TTT that determines how much better a neighbor must be before handover occurs.

\CNdisplay{\text{Effective HO margin} \approx \mathrm{A3\_Offset} + \mathrm{Hysteresis} + \underbrace{\left(\text{UE speed} \times \mathrm{TTT}\right)}_{\text{signal change during TTT}}}

\CNsubsubsection{Example Calculation}

\tcbset{CNcodemode/.style={unbreakable}}\def\CNcodelang{}

\begin{CNplaincode}
\begin{Verbatim}
Given:
  A3 Offset = 3 dB
  Hysteresis = 2 dB  
  TTT = 320 ms
  UE speed = 50 km/h → path loss changes ~0.5 dB during TTT

Effective margin = 3 + 2 + 0.5 = 5.5 dB

Meaning: Neighbor must be 5.5 dB better than serving 
         before handover is triggered
\end{Verbatim}
\end{CNplaincode}

\Needspace{14\baselineskip}

\CNsection{6}{Handover Procedure — LTE X2 Handover}

\CNchained\CNsubsection{}{WHY}

The X2-based handover is the \textbf{preferred} handover mechanism in LTE because it enables direct eNB-to-eNB communication without routing through the core network (MME), reducing latency and signaling load.

\CNsubsection{}{WHAT}

\begin{itemize}
\tightlist
\item
  Flat architecture: eNBs communicate directly via X2 interface
\item
  UE-assisted, network-controlled: UE measures, eNB decides
\item
  Preparation phase ensures target has resources BEFORE handover command
\item
  Total interruption time: \textbf{20–50~ms} (imperceptible for most applications)
\end{itemize}

\CNkeepfig{m04-x2-handover}{3}

\CNsubsection{}{HOW — X2 Handover Sequence}

\CNfigure{m04-x2-handover}{LTE X2 handover: preparation, execution and completion}

\Needspace{11\baselineskip}

\Needspace{19\baselineskip}

\CNsubsection{}{Phase-by-Phase Breakdown}

\CNchained\CNsubsubsection{Phase 1: Preparation (\textasciitilde{}50–200~ms)}

{\def\LTcaptype{none} % do not increment counter
\Needspace{13\baselineskip}
\begin{longtable}[c]{@{}ccc@{}}
\toprule\noalign{}
\rowcolor{CNink}\CNth{Step} & \CNth{Message} & \CNth{Purpose} \\
\CNheadrulesep
\endhead
\bottomrule\noalign{}
\endlastfoot
1 & Measurement Report & UE informs source about neighbor quality \\
2 & HO Decision & Source eNB runs algorithm based on report \\
3 & Handover Request (X2) & Source asks target to prepare resources \\
4 & Admission Control & Target checks if it can accept the UE \\
5 & HO Request Ack (X2) & Target confirms, provides HO command \\
\end{longtable}
}

\Needspace{16\baselineskip}

\CNsubsubsection{Phase 2: Execution (\textasciitilde{}20–50~ms interruption)}

{\def\LTcaptype{none} % do not increment counter
\Needspace{13\baselineskip}
\begin{longtable}[c]{@{}ccc@{}}
\toprule\noalign{}
\rowcolor{CNink}\CNth{Step} & \CNth{Message} & \CNth{Purpose} \\
\CNheadrulesep
\endhead
\bottomrule\noalign{}
\endlastfoot
6 & RRC Reconfiguration & Source commands UE to handover \\
7 & SN Status Transfer & Ensures no data loss (PDCP SN alignment) \\
8 & Data Forwarding & Buffered data forwarded to target \\
9 & Random Access & UE synchronizes to target cell \\
10 & Reconfig Complete & UE confirms successful handover \\
\end{longtable}
}

\Needspace{12\baselineskip}

\CNsubsubsection{Phase 3: Completion (\textasciitilde{}20–100~ms)}

{\def\LTcaptype{none} % do not increment counter
\Needspace{9\baselineskip}
\begin{longtable}[c]{@{}ccc@{}}
\toprule\noalign{}
\rowcolor{CNink}\CNth{Step} & \CNth{Message} & \CNth{Purpose} \\
\CNheadrulesep
\endhead
\bottomrule\noalign{}
\endlastfoot
11 & Path Switch Request & Informs MME/\allowbreak{}S-GW of new location \\
12 & Path Switch Ack & Core confirms path update \\
13 & UE Context Release & Source releases resources \\
\end{longtable}
}

\Needspace{17\baselineskip}

\CNsubsection{}{Key Timing Values}

{\def\LTcaptype{none} % do not increment counter
\Needspace{13\baselineskip}
\begin{longtable}[c]{@{}cc@{}}
\toprule\noalign{}
\rowcolor{CNink}\CNth{Parameter} & \CNth{Typical Value} \\
\CNheadrulesep
\endhead
\bottomrule\noalign{}
\endlastfoot
Total preparation time & 50–200~ms \\
\textbf{User-plane interruption} & \textbf{20–50~ms} \\
Path switch completion & 20–100~ms \\
Total end-to-end & 100–350~ms \\
RACH preamble & Dedicated (no contention) or contention-based \\
\end{longtable}
}

\CNsubsection{}{S1-based Handover (Fallback)}

When X2 interface is NOT available (no direct connection between eNBs):

\begin{itemize}
\tightlist
\item
  All messages route through MME
\item
  Longer preparation time (\textasciitilde{}100–500~ms)
\item
  Used for inter-eNB handover without X2, or inter-MME handover
\item
  Same UE experience but higher network latency
\end{itemize}

\Needspace{14\baselineskip}

\CNsection{7}{Handover Failures}

\CNchained\CNsubsection{}{WHY}

Handover failures directly impact user experience — dropped calls, interrupted video, lost data. Understanding failure modes is essential for network optimization.

\CNsubsection{}{WHAT}

Four primary failure modes exist, each with distinct causes and solutions:

\Needspace{11\baselineskip}

\CNsubsection{7.1}{Too-Late Handover}

\CNchained\CNsubsubsection{WHY it happens}

The UE moves too fast or the handover margin is too large, so the serving cell signal drops below usable level BEFORE the handover completes.

\CNkeepfig{m04-too-late}{2}

\CNsubsubsection{WHAT happens}

\CNfigure{m04-too-late}{Too-late handover: the serving link fails before the handover}

\CNsubsubsection{Indicators}

\begin{itemize}
\tightlist
\item
  Radio Link Failure (RLF) in source cell BEFORE HO command received
\item
  High RLF counter on source cell
\item
  RLF report shows strong neighbor at time of failure
\end{itemize}

\CNsubsubsection{Mitigation}

\begin{itemize}
\tightlist
\item
  \textbf{Decrease} A3 offset (trigger HO earlier)
\item
  \textbf{Decrease} TTT (react faster)
\item
  \textbf{Decrease} hysteresis
\item
  Add early measurement reporting (A2 event with lower threshold)
\end{itemize}

\Needspace{11\baselineskip}

\CNsubsection{7.2}{Too-Early Handover}

\CNchained\CNsubsubsection{WHY it happens}

Handover triggers too quickly (parameters too aggressive). UE is handed to a cell that is not yet strong enough to sustain the connection.

\CNkeepfig{m04-too-early}{2}

\CNsubsubsection{WHAT happens}

\CNfigure{m04-too-early}{Too-early handover: the target cell cannot yet sustain the link}

\CNsubsubsection{Indicators}

\begin{itemize}
\tightlist
\item
  RLF occurs at TARGET cell shortly after handover completion
\item
  UE attempts re-establishment back at source cell
\item
  Short time between HO completion and RLF
\end{itemize}

\CNsubsubsection{Mitigation}

\begin{itemize}
\tightlist
\item
  \textbf{Increase} A3 offset (require bigger difference)
\item
  \textbf{Increase} TTT (wait longer)
\item
  \textbf{Increase} hysteresis
\item
  Verify neighbor cell coverage footprint
\end{itemize}

\Needspace{11\baselineskip}

\CNsubsection{7.3}{Ping-Pong Effect}

\CNchained\CNsubsubsection{WHY it happens}

Two cells have similar signal strength in an overlap area. With aggressive handover parameters, the UE bounces repeatedly between them.

\CNkeepfig{m04-ping-pong}{2}

\CNsubsubsection{WHAT happens}

\CNfigure{m04-ping-pong}{Ping-pong: repeated handovers between two similar cells}

\CNsubsubsection{Impact}

\begin{itemize}
\tightlist
\item
  Excessive signaling load
\item
  Interrupted user data (each HO = 20-50~ms break)
\item
  Wasted radio resources for HO preparation
\item
  Poor QoE for the user
\end{itemize}

\CNsubsubsection{Indicators}

\begin{itemize}
\tightlist
\item
  Same UE performing multiple handovers between same cell pair within short time
\item
  High handover rate with low mobility
\item
  A\ensuremath{\to}B followed by B\ensuremath{\to}A within seconds
\end{itemize}

\Needspace{16\baselineskip}

\CNsubsubsection{Mitigation}

{\def\LTcaptype{none} % do not increment counter
\Needspace{13\baselineskip}
\begin{longtable}[c]{@{}
  >{\raggedright\arraybackslash\hspace{0pt}}m{(\linewidth - 2\tabcolsep) * \real{0.4348}}
  >{\raggedright\arraybackslash\hspace{0pt}}m{(\linewidth - 2\tabcolsep) * \real{0.5652}}@{}}
\toprule\noalign{}
\rowcolor{CNink}\CNth{Strategy} & \CNth{How it helps} \\
\CNheadrulesep
\endhead
\bottomrule\noalign{}
\endlastfoot
\textbf{Increase hysteresis} & Larger margin needed to trigger \ensuremath{\to} fewer triggers \\
\textbf{Increase TTT} & Must sustain longer \ensuremath{\to} transient fluctuations filtered \\
\textbf{Cell-Individual Offset (CIO)} & Bias one cell to be preferred in overlap \\
\textbf{Increase A3 offset} & Neighbor must be significantly better \\
\textbf{Antenna tilt optimization} & Reduce overlap zone size \\
\end{longtable}
}

\textbf{Rule of thumb}: If \textgreater{}30\% of handovers between a cell pair are ping-pong (HO back within 5s), increase TTT or hysteresis for that pair.

\Needspace{11\baselineskip}

\CNsubsection{7.4}{Handover to Wrong Cell}

\CNchained\CNsubsubsection{WHY it happens}

The best target cell is not the one selected. This can occur due to:

\begin{itemize}
\tightlist
\item
  Missing neighbor relation (target not in neighbor list)
\item
  Measurement report doesn’t include the actual best cell
\item
  Propagation anomaly (signal from far cell momentarily strong)
\end{itemize}

\CNsubsubsection{WHAT happens}

\begin{itemize}
\tightlist
\item
  UE is handed to Cell X, but Cell Y would have been the correct target
\item
  UE quickly experiences poor signal at Cell X
\item
  May trigger another handover (to Cell Y) or RLF
\end{itemize}

\CNsubsubsection{Mitigation}

\begin{itemize}
\tightlist
\item
  \textbf{Automatic Neighbor Relation (ANR)}: Automatically discover and add neighbors
\item
  \textbf{Comprehensive neighbor lists}: Ensure all potential targets are configured
\item
  \textbf{Cell-specific offsets}: Deprioritize cells known to cause issues
\item
  \textbf{Layer-3 filtering}: Average measurements to reduce anomalies
\end{itemize}

\CNkeepfig{m04-ho-failure-states}{3}

\CNsubsection{7.5}{Handover Failure State Diagram}

\CNfigure{m04-ho-failure-states}{Handover states and failure modes}

\Needspace{14\baselineskip}

\CNsection{8}{Evolution of Handover: 2G \ensuremath{\to} 3G \ensuremath{\to} LTE \ensuremath{\to} 5G}

\CNchained\CNsubsection{}{WHY}

Each generation improved handover mechanisms to support higher speeds, lower latency, and more complex network architectures. Understanding evolution reveals the design rationale.

\Needspace{11\baselineskip}

\CNkeepfig{m04-gsm-ho-arch}{5}

\CNsubsection{8.1}{GSM (2G): Hard Handover, BSC-Controlled}

\CNchained\CNsubsubsection{Architecture}

\CNfigure{m04-gsm-ho-arch}{GSM: hard handover decided by the BSC}

\Needspace{19\baselineskip}

\CNsubsubsection{Characteristics}

{\def\LTcaptype{none} % do not increment counter
\Needspace{16\baselineskip}
\begin{longtable}[c]{@{}cc@{}}
\toprule\noalign{}
\rowcolor{CNink}\CNth{Aspect} & \CNth{2G GSM} \\
\CNheadrulesep
\endhead
\bottomrule\noalign{}
\endlastfoot
Handover type & \textbf{Hard only} (break-before-make) \\
Decision maker & \textbf{BSC} (Base Station Controller) \\
Measurements & MS measures up to 6 neighbors (BCCH carriers) \\
Measurement rate & Every SACCH period (480~ms) \\
Interruption & 100–200~ms \\
Criteria & RXLEV (signal strength), RXQUAL (BER-based quality) \\
Power control & Combined with handover decision \\
\end{longtable}
}

\CNsubsubsection{How it worked}

\begin{enumerate}
\def\labelenumi{\arabic{enumi}.}
\tightlist
\item
  MS continuously measures serving + 6 strongest neighbors (from BA list)
\item
  MS reports measurements every 480~ms via SACCH
\item
  BSC applies averaging (HREQAVE, HREQT parameters)
\item
  BSC decides based on multiple criteria:

  \begin{itemize}
  \tightlist
  \item
    Power budget (PBGT) \ensuremath{\to} coverage-based
  \item
    RXQUAL \ensuremath{\to} quality-based (interference)
  \item
    RXLEV \ensuremath{\to} level-based (signal too weak/\allowbreak{}strong)
  \item
    Distance \ensuremath{\to} timing advance too large
  \item
    Traffic \ensuremath{\to} congestion relief
  \end{itemize}
\end{enumerate}

\CNsubsubsection{Limitations}

\begin{itemize}
\tightlist
\item
  Hard handover: audible click on voice calls
\item
  BSC-controlled: all intelligence in network, slow adaptation
\item
  Limited measurements: only 6 neighbors, slow reporting
\item
  No soft handover: no diversity gain at cell edge
\end{itemize}

\Needspace{11\baselineskip}

\CNkeepfig{m04-umts-ho-arch}{5}

\CNsubsection{8.2}{UMTS (3G): Soft Handover, RNC-Controlled}

\CNchained\CNsubsubsection{Architecture}

\CNfigure{m04-umts-ho-arch}{UMTS: handover decided by the RNC}

\Needspace{17\baselineskip}

\CNsubsubsection{Characteristics}

{\def\LTcaptype{none} % do not increment counter
\Needspace{14\baselineskip}
\begin{longtable}[c]{@{}cc@{}}
\toprule\noalign{}
\rowcolor{CNink}\CNth{Aspect} & \CNth{3G UMTS} \\
\CNheadrulesep
\endhead
\bottomrule\noalign{}
\endlastfoot
Handover type & \textbf{Soft + Softer + Hard} \\
Decision maker & \textbf{RNC} (Radio Network Controller) \\
Measurements & CPICH Ec/\allowbreak{}No, RSCP \\
Active Set & Up to 3–6 cells simultaneously \\
Interruption & \textbf{0~ms} (soft HO), 20–50~ms (hard HO) \\
Criteria & Events 1A–1C (active set), 2A–2F (inter-freq/\allowbreak{}RAT) \\
\end{longtable}
}

\CNkeepfig{m04-soft-ho-mech}{2}

\CNsubsubsection{Soft Handover Mechanism}

\CNfigure{m04-soft-ho-mech}{Soft handover: two radio legs combined in the RNC}

\begin{itemize}
\tightlist
\item
  \textbf{Active Set Management}:

  \begin{itemize}
  \tightlist
  \item
    Event 1A: $\mathrm{Pilot} > \mathrm{Best\_Pilot} - \mathrm{Reporting\_Range} + \mathrm{Hysteresis}$ \ensuremath{\to} ADD cell
  \item
    Event 1B: $\mathrm{Pilot} < \mathrm{Best\_Pilot} - \mathrm{Reporting\_Range} - \mathrm{Hysteresis}$ \ensuremath{\to} REMOVE cell
  \item
    Event 1C: New pilot \textgreater{} Worst\_\allowbreak{}in\_\allowbreak{}ActiveSet + Hysteresis \ensuremath{\to} REPLACE cell
  \end{itemize}
\item
  \textbf{Macro-diversity gain}: 2–3~dB at cell edge
\item
  \textbf{20–40\%} of connections in soft handover in typical network
\end{itemize}

\CNsubsubsection{Key Innovation}

Soft handover enabled seamless mobility but at the cost of:

\begin{itemize}
\tightlist
\item
  Multiple Channelization Codes consumed (one per active set member)
\item
  Increased Iub backhaul bandwidth
\item
  Complex power control (must satisfy all links)
\end{itemize}

\Needspace{11\baselineskip}

\CNkeepfig{m04-lte-ho-arch}{5}

\CNsubsection{8.3}{LTE (4G): Hard Handover, UE-Assisted Network-Controlled, X2-Based}

\CNchained\CNsubsubsection{Architecture}

\CNfigure{m04-lte-ho-arch}{LTE: flat architecture, the eNB decides and uses X2}

\Needspace{17\baselineskip}

\CNsubsubsection{Characteristics}

{\def\LTcaptype{none} % do not increment counter
\Needspace{14\baselineskip}
\begin{longtable}[c]{@{}cc@{}}
\toprule\noalign{}
\rowcolor{CNink}\CNth{Aspect} & \CNth{4G LTE} \\
\CNheadrulesep
\endhead
\bottomrule\noalign{}
\endlastfoot
Handover type & \textbf{Hard only} (break-before-make) \\
Decision maker & \textbf{Source eNB} (distributed decision) \\
Measurements & RSRP, RSRQ (events A1–A6, B1–B2) \\
Interruption & \textbf{20–50~ms} \\
Interface & \textbf{X2} (direct eNB-to-eNB) \\
Key feature & UE-assisted, network-controlled \\
\end{longtable}
}

\CNsubsubsection{Why LTE Dropped Soft Handover}

\begin{enumerate}
\def\labelenumi{\arabic{enumi}.}
\tightlist
\item
  \textbf{OFDMA} (not CDMA): Different cells use different subcarriers \ensuremath{\to} cannot combine same signal from multiple cells
\item
  \textbf{Flat architecture}: No central RNC to perform combining
\item
  \textbf{Hard HO is fast enough}: 20–50~ms interruption is acceptable for data
\item
  \textbf{Simplicity}: Significantly simpler than managing Active Sets
\item
  \textbf{Efficiency}: No wasted resources for multiple simultaneous links
\end{enumerate}

\CNsubsubsection{Key Innovations in LTE}

\begin{itemize}
\tightlist
\item
  \textbf{X2 interface}: Direct eNB-eNB \ensuremath{\to} fast preparation, data forwarding
\item
  \textbf{Event-based reporting}: Only report when conditions met (less signaling)
\item
  \textbf{Layer-3 filtering}: Configurable averaging to smooth measurements
\item
  \textbf{Automatic Neighbor Relations (ANR)}: Self-configuring neighbor lists
\item
  \textbf{Mobility Robustness Optimization (MRO)}: Self-healing handover parameters
\end{itemize}

\Needspace{11\baselineskip}

\CNkeepfig{m04-nr-ho-arch}{5}

\CNsubsection{8.4}{5G NR: Conditional Handover, DAPS, Dual Connectivity}

\CNchained\CNsubsubsection{Architecture}

\CNfigure{m04-nr-ho-arch}{5G NR: Xn handover and dual connectivity}

\Needspace{16\baselineskip}

\CNsubsubsection{Characteristics}

{\def\LTcaptype{none} % do not increment counter
\Needspace{13\baselineskip}
\begin{longtable}[c]{@{}cc@{}}
\toprule\noalign{}
\rowcolor{CNink}\CNth{Aspect} & \CNth{5G NR} \\
\CNheadrulesep
\endhead
\bottomrule\noalign{}
\endlastfoot
Handover type & Hard, \textbf{Conditional}, DAPS \\
Decision maker & Source gNB (+ CPC for conditional) \\
Measurements & SS-RSRP, SS-RSRQ, SS-SINR, CSI-RS based \\
Interruption & 0~ms (DAPS), \textasciitilde{}0~ms (DC), 20~ms (hard) \\
Key features & CHO, DAPS, Dual Connectivity, MR-DC \\
\end{longtable}
}

\CNkeepfig{m04-cho-compare}{6}

\CNsubsubsection{Conditional Handover (CHO) — 5G Innovation}

\textbf{WHY}: In legacy LTE, the eNB must wait for measurement report, make decision, prepare target, then send command. If radio conditions degrade quickly (high speed, mmWave), this is too slow.

\textbf{WHAT}: Network \textbf{pre-prepares} multiple target cells and gives the UE conditions to trigger the handover \textbf{autonomously}.

\CNfigure{m04-cho-compare}{Traditional versus conditional handover}

\textbf{HOW}:

\begin{enumerate}
\def\labelenumi{\arabic{enumi}.}
\tightlist
\item
  Network configures CHO: ``If Event A3 for Cell X OR Cell Y, execute HO to that cell''
\item
  Target cells are pre-prepared (resources reserved)
\item
  UE monitors conditions
\item
  When condition met \ensuremath{\to} UE immediately executes (no round-trip to network)
\item
  \textbf{Benefit}: Eliminates handover preparation delay \ensuremath{\to} fewer failures at high speed
\end{enumerate}

\CNkeepfig{m04-daps}{6}

\CNsubsubsection{Dual Active Protocol Stack (DAPS) — 0~ms Interruption}

\textbf{WHY}: Even 20~ms interruption matters for URLLC (Ultra-Reliable Low-Latency Communication).

\textbf{WHAT}: UE maintains connection to BOTH source and target simultaneously during handover — similar concept to 3G soft handover but implemented differently.

\CNfigure{m04-daps}{DAPS handover: make before break}

\textbf{HOW}:

\begin{itemize}
\tightlist
\item
  UE has two active protocol stacks (hence the name)
\item
  Data is sent/\allowbreak{}received through both paths during transition
\item
  Once target confirms successful connection, source is released
\item
  \textbf{Trade-off}: Increased UE complexity and power consumption
\end{itemize}

\CNsubsubsection{Dual Connectivity (MR-DC) — Multi-RAT}

\textbf{WHAT}: UE simultaneously connected to two different nodes:

\begin{itemize}
\tightlist
\item
  \textbf{EN-DC} (E-UTRA-NR Dual Connectivity): LTE master + NR secondary
\item
  \textbf{NR-DC}: NR master + NR secondary
\item
  \textbf{NE-DC}: NR master + LTE secondary
\end{itemize}

\CNfigure{m04-dual-connectivity}{Dual connectivity: two radio links at the same time}

\Needspace{20\baselineskip}

\CNsubsection{8.5}{Generation Comparison Summary}

{\def\LTcaptype{none} % do not increment counter
\Needspace{16\baselineskip}
\begin{longtable}[c]{@{}
  >{\raggedright\arraybackslash\hspace{0pt}}m{(\linewidth - 8\tabcolsep) * \real{0.2195}}
  >{\raggedright\arraybackslash\hspace{0pt}}m{(\linewidth - 8\tabcolsep) * \real{0.1951}}
  >{\raggedright\arraybackslash\hspace{0pt}}m{(\linewidth - 8\tabcolsep) * \real{0.2195}}
  >{\raggedright\arraybackslash\hspace{0pt}}m{(\linewidth - 8\tabcolsep) * \real{0.1951}}
  >{\raggedright\arraybackslash\hspace{0pt}}m{(\linewidth - 8\tabcolsep) * \real{0.1707}}@{}}
\toprule\noalign{}
\rowcolor{CNink}\CNth{Feature} & \CNth{2G GSM} & \CNth{3G UMTS} & \CNth{4G LTE} & \CNth{5G NR} \\
\CNheadrulesep
\endhead
\bottomrule\noalign{}
\endlastfoot
HO Type & Hard & Soft + Hard & Hard & Hard + CHO + DAPS \\
Controller & BSC & RNC & eNB & gNB \\
Interruption & 100-200~ms & 0~ms (soft) & 20-50~ms & 0~ms (DAPS/\allowbreak{}DC) \\
Measurement & RXLEV, RXQUAL & CPICH RSCP, Ec/\allowbreak{}No & RSRP, RSRQ & SS-RSRP, CSI-RS \\
Intelligence & Network-only & Network-only & UE-assisted & UE-autonomous (CHO) \\
Architecture & Hierarchical & Hierarchical & Flat & Flat + DC \\
Self-optimization & Manual & Manual & SON (MRO) & AI/\allowbreak{}ML-based \\
\end{longtable}
}

\Needspace{16\baselineskip}

\CNsection{9}{Engineering Examples with Typical Parameter Values}

\CNchained\CNsubsection{9.1}{Example: Urban Macro Network — LTE Handover Optimization}

\CNchained\CNsubsubsection{Scenario}

\begin{itemize}
\tightlist
\item
  Dense urban area, 500m inter-site distance
\item
  Users: mix of pedestrians (3 km/\allowbreak{}h) and vehicles (50 km/\allowbreak{}h)
\item
  Problem: 5\% handover failure rate, excessive ping-pong
\end{itemize}

\Needspace{14\baselineskip}

\CNsubsubsection{Initial Parameters (before optimization)}

{\def\LTcaptype{none} % do not increment counter
\Needspace{11\baselineskip}
\begin{longtable}[c]{@{}ccc@{}}
\toprule\noalign{}
\rowcolor{CNink}\CNth{Parameter} & \CNth{Value} & \CNth{Issue} \\
\CNheadrulesep
\endhead
\bottomrule\noalign{}
\endlastfoot
A3 Offset & 2~dB & Too aggressive \ensuremath{\to} too-early HO \\
Hysteresis & 1~dB & Not enough filtering \\
TTT & 100~ms & Too short \ensuremath{\to} ping-pong \\
Layer-3 filter (k) & 1 & Not enough averaging \\
\end{longtable}
}

\Needspace{16\baselineskip}

\CNsubsubsection{Optimized Parameters}

{\def\LTcaptype{none} % do not increment counter
\Needspace{13\baselineskip}
\begin{longtable}[c]{@{}
  >{\raggedright\arraybackslash\hspace{0pt}}m{(\linewidth - 4\tabcolsep) * \real{0.3333}}
  >{\raggedright\arraybackslash\hspace{0pt}}m{(\linewidth - 4\tabcolsep) * \real{0.3333}}
  >{\raggedright\arraybackslash\hspace{0pt}}m{(\linewidth - 4\tabcolsep) * \real{0.3333}}@{}}
\toprule\noalign{}
\rowcolor{CNink}\CNth{Parameter} & \CNth{New Value} & \CNth{Rationale} \\
\CNheadrulesep
\endhead
\bottomrule\noalign{}
\endlastfoot
A3 Offset & 3~dB & More margin before triggering \\
Hysteresis & 2~dB & Better fading resistance \\
TTT & 256~ms & Eliminates most ping-pong for pedestrians \\
Layer-3 filter (k) & 4 & Smooths fast fading (\textasciitilde{}200~ms averaging) \\
CIO (cell pair with ping-pong) & +3~dB & Biases HO for known problem pair \\
\end{longtable}
}

\CNsubsubsection{Result}

\begin{itemize}
\tightlist
\item
  Handover failure rate: 5\% \ensuremath{\to} 1.5\%
\item
  Ping-pong rate: 12\% \ensuremath{\to} 3\%
\item
  Trade-off: Mean HO delay increased by \textasciitilde{}150~ms (acceptable)
\end{itemize}

\Needspace{11\baselineskip}

\CNsubsection{9.2}{Example: High-Speed Rail — 350 km/\allowbreak{}h}

\CNchained\CNsubsubsection{Challenge}

At 350 km/\allowbreak{}h, the UE crosses a cell (1~km coverage) in \textasciitilde{}10 seconds. Standard parameters are too slow.

\Needspace{16\baselineskip}

\CNsubsubsection{Parameter Tuning}

{\def\LTcaptype{none} % do not increment counter
\Needspace{13\baselineskip}
\begin{longtable}[c]{@{}ccc@{}}
\toprule\noalign{}
\rowcolor{CNink}\CNth{Parameter} & \CNth{Standard} & \CNth{High-Speed} \\
\CNheadrulesep
\endhead
\bottomrule\noalign{}
\endlastfoot
A3 Offset & 3~dB & 2~dB \\
Hysteresis & 2~dB & 1~dB \\
TTT & 256~ms & 40~ms \\
Layer-3 filter (k) & 4 & 1 \\
A2 threshold & -120~dBm & -110~dBm \\
\end{longtable}
}

\CNsubsubsection{Additional Measures}

\begin{itemize}
\tightlist
\item
  Dedicated RF planning along rail corridor
\item
  Directional antennas pointing along track
\item
  Cell overlap zone: 200–300~m minimum
\item
  Consider \textbf{Conditional Handover (5G)} for this scenario
\end{itemize}

\CNsubsubsection{Speed Sensitivity Calculation}

\tcbset{CNcodemode/.style={unbreakable}}\def\CNcodelang{}

\begin{CNplaincode}
\begin{Verbatim}
At 350 km/h = 97 m/s
Path loss change rate ≈ 0.3 dB per 10 ms (at 2 GHz, Hata model)

With TTT = 256 ms: Signal changes by ~7.5 dB during TTT!
→ This is unacceptable — must reduce TTT

With TTT = 40 ms: Signal changes by ~1.2 dB during TTT
→ Acceptable margin
\end{Verbatim}
\end{CNplaincode}

\Needspace{11\baselineskip}

\CNsubsection{9.3}{Example: Inter-RAT Handover — LTE to 3G CSFB}

\CNchained\CNsubsubsection{Scenario}

VoLTE not deployed; voice calls require CS Fallback to 3G.

\Needspace{16\baselineskip}

\CNsubsubsection{Configuration}

{\def\LTcaptype{none} % do not increment counter
\Needspace{13\baselineskip}
\begin{longtable}[c]{@{}
  >{\raggedright\arraybackslash\hspace{0pt}}m{(\linewidth - 4\tabcolsep) * \real{0.4074}}
  >{\raggedright\arraybackslash\hspace{0pt}}m{(\linewidth - 4\tabcolsep) * \real{0.2593}}
  >{\raggedright\arraybackslash\hspace{0pt}}m{(\linewidth - 4\tabcolsep) * \real{0.3333}}@{}}
\toprule\noalign{}
\rowcolor{CNink}\CNth{Parameter} & \CNth{Value} & \CNth{Purpose} \\
\CNheadrulesep
\endhead
\bottomrule\noalign{}
\endlastfoot
B2 Threshold1 (serving LTE) & -110~dBm RSRP & Trigger when LTE weak \\
B2 Threshold2 (UMTS target) & -95~dBm RSCP & Only if 3G is adequate \\
TTT (B2) & 640~ms & Avoid unnecessary fallback \\
A2 Threshold (start measuring 3G) & -105~dBm RSRP & Begin inter-RAT measurements \\
IRAT measurement gap & 6~ms /\allowbreak{} 40~ms & Gap pattern for measuring UMTS \\
\end{longtable}
}

\CNsubsubsection{Procedure}

\begin{enumerate}
\def\labelenumi{\arabic{enumi}.}
\tightlist
\item
  UE on LTE, RSRP drops below -105~dBm \ensuremath{\to} A2 event \ensuremath{\to} start measuring UMTS
\item
  LTE RSRP drops below -110~dBm AND UMTS RSCP \textgreater{} -95~dBm for 640~ms \ensuremath{\to} B2 event
\item
  eNB triggers inter-RAT handover to UMTS
\item
  UE performs hard handover to 3G, establishes CS domain connection
\item
  Voice call proceeds on 3G circuit-switched
\item
  After call ends: UE returns to LTE (fast return /\allowbreak{} redirection)
\end{enumerate}

\Needspace{11\baselineskip}

\CNsubsection{9.4}{Example: Measurement Event Configuration (Complete)}

\CNchained\CNsubsubsection{Full A3 Event Configuration for Intra-Frequency HO}

\tcbset{CNcodemode/.style={unbreakable}}\def\CNcodelang{}

\begin{CNplaincode}
\begin{Verbatim}
RRC Connection Reconfiguration (measurement config):
{
  measObjectEUTRA: {
    carrierFreq: 1850 (EARFCN for Band 3, 1800 MHz)
    allowedMeasBandwidth: 50 RBs
    neighCellConfig: 01 (partial measurement)
    offsetFreq: 0 dB
  }
  
  reportConfigEUTRA: {
    triggerType: event
    eventId: A3
    a3-Offset: 3 dB (= 6 in IE, 0.5 dB steps)
    hysteresis: 2 dB (= 4 in IE, 0.5 dB steps)
    timeToTrigger: 256 ms
    triggerQuantity: RSRP
    reportQuantity: both (RSRP + RSRQ)
    maxReportCells: 4
    reportInterval: 480 ms
    reportAmount: infinity
  }
  
  measIdToAddList: {
    measId: 1
    measObjectId: 1
    reportConfigId: 1
  }
  
  quantityConfig: {
    filterCoefficientRSRP: 4 (k=4, ~200ms averaging)
    filterCoefficientRSRQ: 4
  }
}
\end{Verbatim}
\end{CNplaincode}

\CNsubsubsection{Layer-3 Filtering Formula}

\CNdisplay{F_n = (1 - a)\,F_{n-1} + a\,M_n, \qquad a = \left(\tfrac{1}{2}\right)^{k/4}}

where $k$ is the filterCoefficient (0, 1, 2, …, 8, …):

\begin{itemize}
\tightlist
\item
  $k = 0$: $a = 1.0$ \ensuremath{\to} no filtering (raw measurement used)
\item
  $k = 4$: $a = 0.5$ \ensuremath{\to} 50\% new + 50\% old \ensuremath{\to} about 200~ms memory
\item
  $k = 8$: $a = 0.25$ \ensuremath{\to} 25\% new + 75\% old \ensuremath{\to} about 400~ms memory
\end{itemize}

\Needspace{20\baselineskip}

\CNsubsection{9.5}{Typical Handover KPIs (Operator Benchmarks)}

{\def\LTcaptype{none} % do not increment counter
\Needspace{16\baselineskip}
\begin{longtable}[c]{@{}cccc@{}}
\toprule\noalign{}
\rowcolor{CNink}\CNth{KPI} & \CNth{Target} & \CNth{Good} & \CNth{Poor} \\
\CNheadrulesep
\endhead
\bottomrule\noalign{}
\endlastfoot
Handover Success Rate & \textgreater{}98\% & \textgreater{}99\% & \textless{}95\% \\
Ping-pong Rate & \textless{}5\% & \textless{}2\% & \textgreater{}10\% \\
RLF Rate (HO-related) & \textless{}1\% & \textless{}0.5\% & \textgreater{}3\% \\
Mean Interruption Time & \textless{}50~ms & \textless{}30~ms & \textgreater{}100~ms \\
Too-early HO ratio & \textless{}1\% & \textless{}0.5\% & \textgreater{}3\% \\
Too-late HO ratio & \textless{}2\% & \textless{}1\% & \textgreater{}5\% \\
Unnecessary HO (ping-pong + wrong cell) & \textless{}5\% & \textless{}3\% & \textgreater{}10\% \\
\end{longtable}
}

\Needspace{14\baselineskip}

\CNkeepfig{m04-a3-trigger}{8}

\CNsection{10}{Summary Diagrams}

\CNchained\CNsubsection{}{Measurement Event A3 Triggering — With Hysteresis and TTT}

\CNfigure{m04-a3-trigger}{Event A3 triggering with hysteresis and time-to-trigger}

\CNkeepfig{m04-ho-states}{3}

\CNsubsection{}{Handover States — Complete Flow}

\CNfigure{m04-ho-states}{Complete handover state flow}

\CNkeepfig{m04-param-tree}{3}

\CNsubsection{}{LTE Handover Parameters — Decision Tree}

\CNfigure{m04-param-tree}{Handover parameter optimisation decision tree}

\Needspace{24\baselineskip}

\CNsection{}{Key Formulas Summary}

{\def\LTcaptype{none} % do not increment counter
\Needspace{18\baselineskip}
\begin{longtable}[c]{@{}
  >{\raggedright\arraybackslash\hspace{0pt}}m{(\linewidth - 4\tabcolsep) * \real{0.3600}}
  >{\raggedright\arraybackslash\hspace{0pt}}m{(\linewidth - 4\tabcolsep) * \real{0.4400}}
  >{\raggedright\arraybackslash\hspace{0pt}}m{(\linewidth - 4\tabcolsep) * \real{0.2000}}@{}}
\toprule\noalign{}
\rowcolor{CNink}\CNth{Formula} & \CNth{Expression} & \CNth{Use} \\
\CNheadrulesep
\endhead
\bottomrule\noalign{}
\endlastfoot
A3 Entry & Mn - Hys \textgreater{} Ms + A3\_\allowbreak{}Off & Intra-freq HO trigger \\
A3 Leave & Mn + Hys \textless{} Ms + A3\_\allowbreak{}Off & Cancel pending report \\
Cell Ranking (Rs) & Qmeas,s + QHyst & Serving cell reselection rank \\
Cell Ranking (Rn) & Qmeas,n - Qoffset & Neighbor reselection rank \\
S-criterion & $S_{\mathrm{rxlev}} = Q_{\mathrm{rxlevmeas}} - Q_{\mathrm{rxlevmin}} > 0$ & Cell suitability \\
L3 Filter & $F_n = (1-a)\,F_{n-1} + a\,M_n$ & Measurement smoothing \\
RSRQ & N \ensuremath{\times} RSRP /\allowbreak{} RSSI & Signal quality metric \\
HO Margin & A3\_\allowbreak{}Off + Hys + speed\ensuremath{\times}TTT & Effective trigger margin \\
\end{longtable}
}

\CNboxsection{Review Questions}

\begin{CNbox}[unbreakable]{q}{Review Questions}

\begin{enumerate}
\def\labelenumi{\arabic{enumi}.}
\tightlist
\item
  \textbf{Explain} why LTE uses hard handover instead of soft handover like 3G.
\item
  \textbf{Calculate} the effective handover margin for: A3 offset = 4~dB, Hys = 2~dB, TTT = 320~ms, UE speed = 120 km/\allowbreak{}h.
\item
  \textbf{Draw} the timeline of a too-late handover and identify which parameter to adjust.
\item
  \textbf{Compare} Tracking Areas (LTE) with Location Areas (2G). What is the advantage of Tracking Area Lists?
\item
  \textbf{Design} handover parameters for a highway scenario (130 km/\allowbreak{}h). Justify each choice.
\item
  \textbf{Explain} how Conditional Handover in 5G solves the high-mobility problem that LTE struggles with.
\item
  \textbf{Analyze} a scenario where ping-pong rate is 15\% between two cells. List three mitigation strategies in order of preference.
\item
  \textbf{Describe} the X2 handover procedure and identify where data loss could occur without SN Status Transfer.
\end{enumerate}

\emph{Module 4 complete. Next: Module 5 — Radio Resource Management and Scheduling}

\end{CNbox}
